\chapter{真机策略部署}\label{chap:policy}

本章介绍在真机平台上部署 $\pi_{0.5}$ 策略模型的完整流程，从模型原理出发，到真机环境配置、训练细节，最终给出实验结果与分析。

\section{$\pi_{0.5}$ 模型原理}

$\pi_{0.5}$ \citep{physicalintelligence2025pi05} 是 Physical Intelligence 推出的视觉-语言-动作（Vision-Language-Action, VLA）模型，其核心目标是将视觉感知、语言理解和机器人动作生成统一到一个端到端的网络架构中，实现开放世界场景下的通用操作。

与 RT-2 \citep{zitkovich2023rt2} 等将动作离散化为 token 进行预测的路线不同，$\pi_{0.5}$ 采用 Flow Matching（流匹配）作为动作生成的核心机制。Flow Matching 的基本思想是学习一条从简单先验分布（如标准高斯分布）到目标动作分布之间的连续概率路径。具体而言，在训练阶段，给定一条真实的动作序列 $\mathbf{a} \in \mathbb{R}^{H \times d}$（其中 $H$ 为预测时域，$d$ 为动作维度），模型学习一个时变向量场 $v_\theta(\mathbf{x}_t, t \mid \mathbf{o})$，该向量场定义了从噪声 $\mathbf{x}_0 \sim \mathcal{N}(0, I)$ 到目标动作 $\mathbf{x}_1 = \mathbf{a}$ 的最优传输路径：
\begin{equation}
    \mathbf{x}_t = (1-t)\mathbf{x}_0 + t\mathbf{x}_1, \quad t \in [0,1]
\end{equation}
模型以条件观测 $\mathbf{o}$（包含视觉、语言和本体状态）为条件，预测速度场 $v_\theta$，训练目标为：
\begin{equation}
    \mathcal{L} = \mathbb{E}_{t, \mathbf{x}_0, \mathbf{x}_1} \left[ \| v_\theta(\mathbf{x}_t, t \mid \mathbf{o}) - (\mathbf{x}_1 - \mathbf{x}_0) \|^2 \right]
\end{equation}
推理时，从噪声采样 $\mathbf{x}_0$，通过数值积分（如 Euler 方法）沿学习到的向量场逐步演化，即可生成目标动作序列。相比扩散模型，Flow Matching 所需的推理步数更少（通常 5--10 步即可收敛），有利于满足机器人控制的实时性要求。

本文对 $\pi_{0.5}$ 基座模型进行微调，并使用其官方基座权重作为预训练初始化。该模型本身并不预设动作一定是末端位姿、关节位置或关节增量，而是预测一个连续动作块；在本实现中，模型动作维度为 32，动作预测时域为 50。具体动作语义由机器人数据适配层决定：本文平台的 27 维关节动作被映射到前 27 个动作维度，剩余维度通过零填充与预训练模型的 32 维动作接口对齐。

在多模态输入方面，$\pi_{0.5}$ 的架构由三个核心组件构成。首先是 视觉编码器，采用预训练的 ViT/SigLIP 等视觉 backbone 对多视角图像进行编码。模型使用固定的图像槽位名称（如 \texttt{base\_0\_rgb}、\texttt{right\_wrist\_0\_rgb}、\texttt{left\_wrist\_0\_rgb}）来区分不同视角的输入，并通过图像掩码标记各视角是否有效。其次是 语言 backbone，基于预训练的大语言模型，将自然语言任务指令映射为语义表征，为模型提供任务目标和高层常识先验。最后是 动作头，以 Flow Matching 方式从融合了视觉、语言和本体状态的特征中预测连续动作块。

$\pi_{0.5}$ 的一个关键创新在于其训练策略。除了在多机器人平台的大规模示教数据集上进行预训练外，$\pi_{0.5}$ 在预训练阶段还混入了 大规模网页视觉-语言数据（如 ImageNet、WebLI 等）。这一设计赋予了模型开放世界泛化能力——即使面对训练中未见过的新物体和新场景，模型也能借助语言和视觉先验产生合理的操作行为。此外，$\pi_{0.5}$ 引入了 语义子任务预测（semantic sub-task prediction）作为辅助任务：模型在预测动作的同时，还需判断当前处于任务的哪个语义子阶段。这一设计有效增强了模型对长程任务结构的理解能力，提升了多步操作任务中的连贯性和成功率。

总体来看，$\pi_{0.5}$ 遵循预训练-微调的训练范式：首先在大规模、多 embodiment 的数据集上预训练获得通用操作先验，然后在目标平台的专家示教数据上进行领域微调，使模型适应特定硬件平台和任务场景。

\section{真机实验配置}

我们在我们的实验平台上微调了 $\pi_{0.5}$ 策略。当前采用单臂灵巧操作配置，具体参数详见附录~\ref{app:config}，任务是将紫色圆柱体放入碗内。

\begin{figure}[htb]
    \centering
    \includegraphics[width=0.7\textwidth]{Img/real_policy.jpg}
    \caption{真机实验环境}\label{fig:real_policy}
\end{figure}

在动作空间方面，本文平台采用关节空间作为动作表示，而非末端位姿。需要注意的是，这并不是 $\pi_{0.5}$ 模型本身固定的动作定义，而是我们在 Tianji/WUJI HAND 数据适配层中的选择。具体而言，状态向量与最终输出动作均为 27 维：前 7 维对应右臂关节，后 20 维对应 WUJI HAND 关节。训练时，原始数据中的关节目标为绝对关节位置；为了适配模型训练，代码中将右臂前 7 维转换为相对当前状态的关节增量，而 20 维灵巧手动作保持为绝对关节位置。推理时，模型输出先经过反归一化，再将右臂增量加回当前关节状态，最终发送给控制器的仍是绝对关节目标。采用关节空间的主要原因是灵巧手具有 20 个自由度，若仅控制指尖末端位姿，会丢失手指中间关节的构型信息，容易导致抓取和操作过程中的自碰撞或非自然构型；而全关节表示能够完整保留灵巧操作姿态，并与 WUJI HAND 的底层位置控制接口直接对接。

在视觉输入方面，我们使用两路 RGB 图像：顶部安装的 Intel RealSense D435i 提供俯视全局视角，\texttt{cam\_left\_wrist} 对应的 Orbbec Gemini 336L 提供近距离手部-物体交互视角。当前训练配置中丢弃 \texttt{cam\_right\_wrist}，并将腕部视角映射到 $\pi_{0.5}$ 的 \texttt{right\_wrist\_0\_rgb} 图像槽位；未使用的 \texttt{left\_wrist\_0\_rgb} 槽位以零填充并标记为无效。视觉图像统一缩放至模型要求的输入分辨率（$224 \times 224$），以 \texttt{rgb8} 格式输入。

在数据方面，我们使用第三章所述的数据采集平台获取遥操作专家示教数据，并构建 \texttt{right\_pick} 数据集用于微调。该数据集共包含 49 条有效示教轨迹，总时长约 425.0 s；单条轨迹时长为 4.75--11.83 s，平均时长约 8.67 s。原始数据包含顶部相机、左腕相机和右腕相机三路 RGB 图像，以及右臂关节状态、右手关节状态和遥操作控制命令等 ROS 话题。训练时使用顶部相机和腕部相机两路图像，并将同步后的右臂 7 维关节与右手 20 维关节组成 27 维本体状态和动作序列，再按照 $\pi_{0.5}$ 的输入规范进行时间对齐、归一化和动作块构造。

\section{策略训练}

策略训练使用 8 张 NVIDIA H20\mbox{-}3e GPU 完成全量微调，训练过程中显存占用约 13 GiB/卡，总训练耗时约 6 小时。更详细的硬件环境与超参数配置见附录~\ref{app:config}。

\section{实验结果与分析}

微调后的 $\pi_{0.5}$ 策略已经能够在部分真机试次中完成将紫色圆柱体放入碗内的任务，但整体成功率仍然较低。实验中观察到的失败现象主要集中在以下三类：
\begin{enumerate}
    \item 抓取目标定位错误。策略有时会在圆柱体之外的位置闭合灵巧手，或者抓取圆柱体的错误部位，导致后续无法稳定搬运。
    \item 手部高度过低。执行过程中手部轨迹有时低于期望高度，造成灵巧手或被抓物体与桌面发生碰撞。
    \item 动作输出退化。部分试次中机械臂和灵巧手几乎不运动，策略没有生成有效的接近、抓取或放置动作。
\end{enumerate}

从现有实验现象看，失败原因除了数据量较少以外，主要来自数据质量和视觉观测两方面。首先，训练数据本身可能存在系统性偏差。数据采集与执行过程中机械臂工作在阻抗控制模式下，若末端执行器的重力补偿或负载补偿不充分，实际末端位置会受到灵巧手和末端结构质量的影响而低于期望轨迹。这样得到的示教数据会把偏低的手部姿态带入训练集，使策略在真机执行时更容易复现靠近桌面的轨迹，从而增加碰撞概率。其次，实验现场光照条件较差，顶部相机和腕部相机获得的图像对比度不足，圆柱体和桌面背景之间的视觉区分度降低。这会影响模型对目标位置和抓取部位的判断，进而表现为抓取空位置、抓取圆柱体错误部位，或在视觉输入偏离训练分布时输出接近静止的动作。

因此，后续改进应优先从两方面展开：一方面重新检查阻抗控制和末端负载补偿，减少示教轨迹与实际执行轨迹之间的高度偏差；另一方面改善实验光照并扩充不同光照条件下的数据，以提高视觉输入的稳定性和策略对场景变化的鲁棒性。
